\documentclass[11pt]{article}
\usepackage[T1]{fontenc}
\usepackage[utf8]{inputenc}
\usepackage[margin=1in]{geometry}
\usepackage{amsmath,amssymb}
\usepackage{microtype}
\usepackage[hidelinks]{hyperref}
\newcommand{\Pb}{\mathbb P}
\DeclareMathOperator{\Beta}{Beta}
\DeclareMathOperator{\BR}{BR}
\allowdisplaybreaks[1]
\title{Two-Seller Competition: Oblivious equilibrium}
\author{Progress note --- transcription of handwritten notes}
\date{\today}

\begin{document}
\maketitle
In this note we present four different models of competition where the sellers don't observe the full state. We restrict this note to the forgetting case. 

\section{Elements of the model}

\begin{itemize}
\item Two symmetric providers, indexed by $j\in\{A,B\}$.
\item Two products, indexed by $i\in\{1,2\}$.
\item Success probabilities: $0<p_1<p_2<1$.
\item Costs: $c_1<c_2$.
\item Revenue: $R>c_2$.
\end{itemize}

\section{Customer beliefs.}
Recall that in the forgetting model, past observations are discounted by a factor $\rho$ whenever a new outcome is observed. Thus, after $k$ interactions with the seller, the total effective information satisfies
$$
S_k+F_k=\rho(S_{k-1}+F_{k-1})+1,
$$
with \(S_0+F_0=0\). Iterating this expression yields
$$
S_k+F_k=1+\rho+\cdots+\rho^{k-1}
=\frac{1-\rho^k}{1-\rho} \quad \to 1/(1-\rho)\quad\text{as }k\to\infty.
$$

To study the limiting behavior, we normalize the state as 
$$x_s=(1-\rho)S_k$$ 
$$x_f=(1-\rho)F_k$$

Since $x_s+x_f=1-\rho^k\to1$, we focus on the limiting boundary $x_s+x_f=1$, which can be represented by a single variable $x=x_s$, with $x_f=1-x$. Then $x\in [0,1]$ when $\rho \in [1/2,1)$. 

For each seller we will call this normalized state variable $x_j$ for $j\in\{A,B\}$ that will represent the proportion of successes for provider $j$ in the long run. The customer observes the state
$$
x=(x_A,x_B),
$$
Her beliefs are
\begin{align}
\theta_A\mid x_A&\sim\Beta(1+\frac{x_A}{1-\rho},1+\frac{1-x_A}{1-\rho}),\\
\theta_B\mid x_B&\sim\Beta(1+\frac{x_B}{1-\rho},1+\frac{1-x_B}{1-\rho}).
\end{align}


\paragraph{Customer choice.}
The customer uses Thompson sampling (TS), drawing independently from the
two belief distributions and choosing
\begin{equation}
y_t=
\begin{cases}
A,&\widetilde\theta_A>\widetilde\theta_B,\\
B,&\widetilde\theta_B>\widetilde\theta_A.
\end{cases}
\end{equation}
Thus the selection probabilities are
\begin{equation}
D_A(x_A,x_B)=\Pb(\widetilde\theta_A>\widetilde\theta_B\mid x_A,x_B),
\qquad D_B(x_A,x_B)=1-D_A(x_A,x_B).
\end{equation}

\section{Timing, state transitions, and strategies}

\paragraph{Timing.}
\begin{enumerate}
    \item The customer observes $x=(x_A,x_B)$, while each
    seller $j$ observes only its own reputation $x_j$.
    
    \item The customer chooses a seller using Thompson sampling.
    
    \item The selected seller chooses a product $i\in\{1,2\}$
    according to its strategy $\sigma_j(\cdot\mid x_j)$
    and receives profit $R-c_i$.
    
    \item The customer observes the outcome, but not the product.
    
    \item The selected seller's reputation is updated, while
    the other seller's reputation remains unchanged.
\end{enumerate}

\paragraph{State transitions.}
Superscripts $s$ and $f$ denote success and failure,
respectively. The state transitions are
\begin{equation}
\begin{aligned}
T_A^s(x)&=(\rho x_A+1-\rho,x_B),\\
T_A^f(x)&=(\rho x_A,x_B),\\
T_B^s(x)&=(x_A,\rho x_B+1-\rho),\\
T_B^f(x)&=(x_A,\rho x_B).
\end{aligned}
\end{equation}

\paragraph{Strategies.}
We consider stationary mixed strategies that depend only
on the seller's own reputation. For each seller $j\in\{A,B\}$,
a strategy is a mapping
\begin{equation}
\sigma_j:[0,1]\rightarrow\Delta(\{1,2\}),
\end{equation}
where $\sigma_j(i\mid x_j)$ denotes the probability that
seller $j$ chooses product $i$ when selected at reputation
$x_j$.

\section{Oblivious equilibrium: Distribution}

\paragraph{Oblivious approximation:}
Suppose that each seller $j\in\{A,B\}$ assumes that
its competitor's reputation follows a stationary distribution
$\mu_{-j}$ on $[0,1]$.

Instead of observing the competitor's current reputation,
seller $j$ computes its expected demand by averaging over
this distribution:
\begin{equation}
d_j(x_j;\mu_{-j})
=
\int_0^1 D(x_j,y)\,\mu_{-j}(dy).
\end{equation}

Thus, each seller's demand depends only on its own reputation
and the conjectured distribution of its competitor.
When optimizing, the seller treats this distribution as fixed.

\paragraph{Seller's problem:}
Let $U_j(x_j)$ denote seller $j$'s optimal expected discounted
profit before customer selection. Given the conjectured distribution $\mu_{-j}$, the Bellman equation is
\begin{equation}
\begin{aligned}
U_j(x_j) ={}& \gamma[1-d_j(x_j;\mu_{-j})]U_j(x_j)\\ &+ d_j(x_j;\mu_{-j})
\max_{i\in\{1,2\}}\Big\{ R-c_i\\
&\qquad+\gamma\left[ p_i U_j(T_\rho^s(x_j)) + (1-p_i)U_j(T_\rho^f(x_j))
\right] \Big\}.
\end{aligned}
\end{equation}

If the seller is not selected, it receives no profit and its reputation remains unchanged. If selected, it chooses a product to maximize its current profit and expected future value.

\paragraph{Stationary distribution:}
Given a strategy $\sigma_j$ and a conjectured distribution
$\mu_{-j}$, the probability of success conditional on
seller $j$ being selected is
\begin{equation}
p_{\sigma_j}(x_j) = \sigma_j(1\mid x_j)p_1 + \sigma_j(2\mid x_j)p_2.
\end{equation}

Given a strategy profile $\sigma=(\sigma_A,\sigma_B)$,
let $P^\sigma(x,C)$ denote the probability that the
joint reputation state belongs to a measurable set
$C\subseteq[0,1]^2$ in the next period, given the
current state $x=(x_A,x_B)$. The transition probabilities are
\begin{equation}
\begin{aligned}
P^\sigma(x,C)
={}&D(x_A,x_B)
\Big[
p_{\sigma_A}(x_A)\mathbf{1}\{T_A^s(x)\in C\} + (1-p_{\sigma_A}(x_A))
\mathbf{1}\{T_A^f(x)\in C\}\Big]\\
&+[1-D(x_A,x_B)]\Big[p_{\sigma_B}(x_B)\mathbf{1}\{T_B^s(x)\in C\}+(1-p_{\sigma_B}(x_B))\mathbf{1}\{T_B^f(x)\in C\}
\Big]
\end{aligned}
\end{equation}

\paragraph{Oblivious equilibrium:}
An oblivious equilibrium is a mixed strategy $\sigma^*$,
a reputation distribution $\mu^*$, and a joint stationary
distribution $\nu^*$ satisfying the following conditions.

\textit{Optimality:}
\begin{equation}
\begin{aligned}
\operatorname{supp}\sigma^*(\cdot\mid x)
\subseteq
\arg\max_{i\in\{1,2\}}
\Big\{
R-c_i
+\gamma\big[
&p_iU_{\mu^*}(T_\rho^s(x))\\
&+(1-p_i)U_{\mu^*}(T_\rho^f(x))
\big]
\Big\}.
\end{aligned}
\end{equation}

\textit{Consistency:}
\begin{equation}
\begin{aligned}
\nu^*&=\nu^*P^{\sigma^*},\\
\mu^*(C)&=\nu^*(C\times[0,1])
=\nu^*([0,1]\times C),
\qquad \forall C\subseteq[0,1]
\text{ measurable}.
\end{aligned}
\end{equation}

Here $P^{\sigma^*}$ denotes the joint transition kernel
of the actual two-seller game when both sellers follow
$\sigma^*$.

The first condition requires each seller to choose only
optimal products given the conjectured distribution $\mu^*$.
The second requires this distribution to coincide with the
marginal distributions generated when both sellers follow
$\sigma^*$ in the actual game.

\section{Oblivious equilibrium: Fixed quality}

\paragraph{Seller's conjecture:}
Suppose that each seller $j\in\{A,B\}$ represents its
competitor by a fixed probability $p_{0,j}\in(0,1)$,
corresponding to the long-run average of the competitor's
Thompson sampling draws.

Instead of considering the competitor's reputation
distribution, seller $j$ computes its expected demand as
\begin{equation}
d(x_j;p_{0,j}) = \mathbb{P}(\widetilde{\theta}_j>p_{0,j}\mid x_j).
\end{equation}

Thus, each seller faces the original problem with an
outside option of fixed quality $p_{0,j}$. When optimizing,
the seller treats this probability as fixed.


\paragraph{Seller's problem:}
Let $U_{p_{0,j}}(x_j)$ denote seller $j$'s optimal expected
discounted profit before customer selection, given
$p_{0,j}$.

The Bellman equation is
\begin{equation}
\begin{aligned}
U_{p_{0,j}}(x_j) ={}& \gamma[1-d(x_j;p_{0,j})]U_{p_{0,j}}(x_j)\\
& + d(x_j;p_{0,j}) \max_{i\in\{1,2\}} \Big\{ R-c_i\\
&\qquad+ \gamma\left[ p_iU_{p_{0,j}}(T_\rho^s(x_j)) + (1-p_i)U_{p_{0,j}}(T_\rho^f(x_j)) \right] \Big\}.
\end{aligned}
\end{equation}


\paragraph{Oblivious equilibrium:}
We consider a symmetric equilibrium in which both sellers
use the same mixed strategy $\sigma^*$ and the same
conjectured probability $p_0^*$. An oblivious equilibrium is a triple
$(\sigma^*,p_0^*,\nu^*)$ satisfying the following conditions.

\textit{Optimality.}
\begin{equation}
\begin{aligned}
\operatorname{supp}\sigma^*(\cdot\mid x)
\subseteq
\arg\max_{i\in\{1,2\}}
\Big\{
R-c_i
+\gamma\big[
&p_iU_{p_0^*}(T_\rho^s(x))\\
&+(1-p_i)U_{p_0^*}(T_\rho^f(x))
\big]
\Big\}.
\end{aligned}
\end{equation}

\textit{Consistency.}
\begin{equation}
\begin{aligned}
\nu^*&=\nu^*P^{\sigma^*},\\
p_0^*
&=
\int_{[0,1]^2}
\frac{1-\rho+x_A}{3-2\rho}
\,\nu^*(dx_A,dx_B)\\
&=
\int_{[0,1]^2}
\frac{1-\rho+x_B}{3-2\rho}
\,\nu^*(dx_A,dx_B).
\end{aligned}
\end{equation}

Here $P^{\sigma^*}$ denotes the joint transition kernel
of the actual two-seller game when both sellers follow
$\sigma^*$.

The optimality condition requires each seller to assign positive probability only to products that maximize its expected discounted profit given $p_0^*$. The consistency condition requires the joint reputation distribution to be stationary and $p_0^*$ to coincide with the long-run average of the competitor's Thompson sampling draws.



\section{Oblivious equilibrium: Average success probability}

\paragraph{Seller's conjecture:}
Suppose that each seller $j\in\{A,B\}$ represents its
competitor by a fixed probability $p_{0,j}\in[p_1,p_2]$,
corresponding to the competitor's long-run success rate
among its sales.

Instead of considering the competitor's reputation
distribution, seller $j$ computes its expected demand as
\begin{equation}
d(x_j;p_{0,j}) = \mathbb{P}(\widetilde{\theta}_j>p_{0,j}\mid x_j).
\end{equation}

Thus, each seller faces the original forgetting monopoly
problem with an outside option of fixed quality $p_{0,j}$.
When optimizing, the seller treats this probability as fixed. Unlike the
previous model, $p_{0,j}$ represents the competitor's
average probability of success among its sales, rather
than the average of its Thompson sampling draws.

\paragraph{Seller's problem:}
Let $U_{p_{0,j}}(x_j)$ denote seller $j$'s optimal expected
discounted profit before customer selection, given
$p_{0,j}$.

The Bellman equation is
\begin{equation}
\begin{aligned}
U_{p_{0,j}}(x_j) ={}& \gamma[1-d(x_j;p_{0,j})]U_{p_{0,j}}(x_j)\\
& + d(x_j;p_{0,j}) \max_{i\in\{1,2\}} \Big\{ R-c_i\\
&\qquad+ \gamma\left[ p_iU_{p_{0,j}}(T_\rho^s(x_j)) + (1-p_i)U_{p_{0,j}}(T_\rho^f(x_j)) \right] \Big\}.
\end{aligned}
\end{equation}

The seller solves the same problem as in the fixed-quality
model. The difference is the consistency condition used
to determine $p_{0,j}$.

\paragraph{Oblivious equilibrium:}
We consider a symmetric equilibrium in which both sellers
use the same mixed strategy $\sigma^*$ and the same
conjectured probability $p_0^*$. An oblivious equilibrium is a triple
$(\sigma^*,p_0^*,\nu^*)$ satisfying the following conditions.

\textit{Optimality.}
\begin{equation}
\begin{aligned}
\operatorname{supp}\sigma^*(\cdot\mid x)
\subseteq
\arg\max_{i\in\{1,2\}}
\Big\{
R-c_i
+\gamma\big[
&p_iU_{p_0^*}(T_\rho^s(x))\\
&+(1-p_i)U_{p_0^*}(T_\rho^f(x))
\big]
\Big\}.
\end{aligned}
\end{equation}

\textit{Consistency.}
\begin{equation}
\begin{aligned}
\nu^*&=\nu^*P^{\sigma^*},\\
p_0^*
&=
\frac{\displaystyle\int_{[0,1]^2}
D_B(x_A,x_B)\,p_{\sigma^*}(x_B)\,\nu^*(dx_A,dx_B)}
{\displaystyle\int_{[0,1]^2}
D_B(x_A,x_B)\,\nu^*(dx_A,dx_B)}\\
&=
\frac{\displaystyle\int_{[0,1]^2}
D_A(x_A,x_B)\,p_{\sigma^*}(x_A)\,\nu^*(dx_A,dx_B)}
{\displaystyle\int_{[0,1]^2}
D_A(x_A,x_B)\,\nu^*(dx_A,dx_B)}.
\end{aligned}
\end{equation}

Here $P^{\sigma^*}$ denotes the joint transition kernel
of the actual two-seller game when both sellers follow
$\sigma^*$, and $p_{\sigma^*}(x)$ is the probability of
success conditional on being selected at reputation $x$.

The optimality condition requires each seller to assign
positive probability only to products that maximize its
expected discounted profit given $p_0^*$. The consistency
condition requires the joint reputation distribution to
be stationary and $p_0^*$ to coincide with the competitor's
success rate among its sales in the actual game. Weighting
by the selection probability and dividing by the seller's
purchase frequency makes this an average among sales.
The probability $p_0^*$ is generated by the strategy and
does not represent an intrinsic quality of the seller.

\section{Oblivious equilibrium: Conditional distribution}

\paragraph{Seller's conjecture:}
Suppose that each seller $j\in\{A,B\}$ assumes that
its competitor's reputation follows a conditional distribution
$\kappa_j(\cdot\mid x_j)$ on $[0,1]$, indexed by the
seller's own reputation $x_j$ before customer selection.

Seller $j$ computes its expected demand as
\begin{equation}
d_j(x_j;\kappa_j)
=
\int_0^1 D(x_j,y)\,\kappa_j(dy\mid x_j).
\end{equation}

The seller uses its own reputation as information about
its competitor, without observing the competitor's current
state. When optimizing, it holds the entire kernel $\kappa_j$
fixed, using its prescribed conditional law at each future
own state.

\paragraph{Seller's problem:}
Let $U_{\kappa_j}(x_j)$ denote seller $j$'s optimal expected
discounted profit before customer selection, given $\kappa_j$.
The Bellman equation is
\begin{equation}
\begin{aligned}
U_{\kappa_j}(x_j) ={}& \gamma[1-d_j(x_j;\kappa_j)]U_{\kappa_j}(x_j)\\
&+d_j(x_j;\kappa_j)\max_{i\in\{1,2\}}\Big\{R-c_i\\
&\qquad+\gamma\left[p_iU_{\kappa_j}(T_\rho^s(x_j))
+(1-p_i)U_{\kappa_j}(T_\rho^f(x_j))\right]\Big\}.
\end{aligned}
\end{equation}

\paragraph{Oblivious equilibrium:}
We consider a symmetric equilibrium in which both sellers
use the same mixed strategy $\sigma^*$ and the same
conjectured conditional kernel $\kappa^*$.
A conditional oblivious equilibrium is a triple
$(\sigma^*,\kappa^*,\nu^*)$, where $\nu^*$ is a symmetric
joint reputation distribution with common marginal $\mu^*$,
satisfying the following conditions.

\textit{Optimality.}
\begin{equation}
\begin{aligned}
\operatorname{supp}\sigma^*(\cdot\mid x)
\subseteq
\arg\max_{i\in\{1,2\}}
\Big\{
R-c_i
+\gamma\big[
&p_iU_{\kappa^*}(T_\rho^s(x))\\
&+(1-p_i)U_{\kappa^*}(T_\rho^f(x))
\big]
\Big\}.
\end{aligned}
\end{equation}

\textit{Consistency.}
\begin{equation}
\begin{aligned}
\nu^*&=\nu^*P^{\sigma^*},\\
\mu^*(C)&=\nu^*(C\times[0,1])
=\nu^*([0,1]\times C),\\
\nu^*(C\times E)
&=\int_C\kappa^*(E\mid x)\,\mu^*(dx),
\qquad C,E\subseteq[0,1]\text{ measurable}.
\end{aligned}
\end{equation}

Here $P^{\sigma^*}$ denotes the joint transition kernel
of the actual two-seller game when both sellers follow
$\sigma^*$. The last condition requires the conjectured
kernel to coincide with the conditional distribution
generated by this joint stationary law, $\mu^*$-almost
everywhere.

On a finite reputation grid, this means
$\kappa^*(y\mid x)=\nu^*(x,y)/\mu^*(x)$ whenever $\mu^*(x)>0$.
At zero-mass grid states, we complete the conjecture using
the unconditional marginal $\mu^*$ and require optimality
under that completion.

The optimality condition requires each seller to choose
only optimal products given $\kappa^*$. Unlike the distribution
model, consistency is imposed on the conditional law for
$\mu^*$-almost every own state, rather than only on the
unconditional marginal. The policy still depends only on
own reputation; the seller does not observe the full joint
state or track a belief based on the entire history of play.

\end{document}
